\documentclass[10pt]{article}
\usepackage[T1]{fontenc}
\usepackage[utf8]{inputenc}
\usepackage{lmodern,amsmath,amssymb}
\usepackage[margin=19mm]{geometry}
\usepackage{xurl}
\usepackage[colorlinks=true,urlcolor=blue]{hyperref}
\usepackage{microtype}
\setlength{\parindent}{0pt}
\setlength{\parskip}{5pt}
\pagestyle{empty}
\hypersetup{pdftitle={RHO5: note for expert review},pdfauthor={Qianli Ma, Chao Wu}}
\begin{document}
{\Large\bfseries RHO5: note for expert review}\par
Qianli Ma\textsuperscript{1,2}, Chao Wu\textsuperscript{1}\quad---\quad11 September 2026\\
\textsuperscript{1}Zhejiang University\qquad\textsuperscript{2}WUJIE AI\\
\textit{Email address (Qianli Ma):} \href{mailto:qianli.ma@zju.edu.cn}{\texttt{qianli.ma@zju.edu.cn}}

\textbf{Purpose.} We would appreciate a critical reading of the attached
manuscript and supplementary index, particularly the global reduction and the interfaces
between analytic lemmas and exact computational certificates.

\textbf{Claim.} For real \(5\times5\) matrices, exact Gaussian elimination
with complete pivoting has maximum element growth
\[
\rho_5^{\mathbb R}=\alpha
=4.132517078632472854223346853277\ldots.
\]
The upper bound quantifies over every legal pivot path, including ties
and singular termination. The constant is the unique root in \((4,5)\)
of the degree-61 polynomial identified by Chen, Edelman, and Urschel.
Their attained lower bound is prior work; the manuscript supplies the
global upper bound and hence the equality in their Conjecture~2.1
(arXiv version~2 numbering).

\textbf{Proof route and reading order.}
\begin{enumerate}
\setlength{\itemsep}{2pt}
\item \textit{Sections 3--5: analytic global entry.}
After max-entry normalization, the first four pivot magnitudes are
at most~4. Actual diagonal contractions
preserve a nonzero fifth pivot and all earlier complete-pivoting
conditions, producing one of two saturated tail geometries, \(D\) or \(X\).
Zero and tied boundary cases are treated explicitly.
\item \textit{Sections 6--7: the two remaining domains.}
The complete \(X\) bound is assembled from identified exact components.
For each nonempty high-value fibre in the normalized range \(p\ge1\),
the remaining negative \(D\) domain admits an attained maximum over six
variables at a fixed seventeen-variable framework; no
historical rank classification is an extra premise of this reduction.
\item \textit{Sections 8--10: finite acceptance and assembly.}
Exact certificates establish either infeasibility of
\(F\ge\gamma=4132517/10^6<\alpha\), or safety at \(\alpha\).
Verified source bindings compose their domains into a complete cover.
The accepted \(X\) and negative-\(D\) bounds then apply to every original
matrix and combine with the attained lower bound.
\end{enumerate}

\textbf{Questions on which feedback would be especially useful.}
\begin{itemize}
\setlength{\itemsep}{2pt}
\item Do the analytic reductions retain the full quantifiers,
including zero coordinates, ties, and common-matrix feasibility?
\item Are the certificate rules and their stated source domains
sufficiently explicit to support the global inference?
What additional presentation or independent checking would you require?
\item Is the claimed advance over the existing fifth-pivot result
correctly positioned, and are important references or qualifications missing?
\end{itemize}

\textbf{Evidence and present scope.} The manuscript and supplementary index identify the exact
polynomial, model, accepted components, and final coverage records.
Historical exact acceptances are reused with explicit identities;
the final composition is not a fresh replay of every earlier tree.
The final \(B\) archives have been preserved locally and hash-checked;
the full theorem also requires the identified upstream \(X\) payloads.
Source, certificate downloads and step-by-step instructions are public at
\url{https://github.com/hkjtsgmc79-boop/rho5-proof} (release v1.0.0).
The guide states the remaining full-chain portability limitations.
The material is supplied for expert feedback. A companion Lean development
formalizes key analytic results and a conditional final theorem; full
global safety still relies on the written arguments and exact certificates.

\textbf{Related paper.} Chen--Edelman--Urschel,
\textit{IMA Journal of Numerical Analysis} (2026), advance article drag039,
published online 7 September 2026.\par
\url{https://doi.org/10.1093/imanum/drag039}\quad
Open manuscript: \url{https://arxiv.org/abs/2602.20390v2}.
\end{document}
